\section{Правила остановки}

\subsection{Общий вид}

Предсказанное распределение \(\widehat p_{t,k}\) само по себе не является правилом действия. Правило остановки --- это функция, которая на шаге \(k\) принимает решение по \(\Hh_{t,k}\) и \(\widehat p_{t,k}\). Реализуемое время остановки имеет вид
\begin{equation}
  \tau_t=\inf\{k:\varphi_k(\widehat p_{t,k},\Hh_{t,k})=1\}.
  \label{eq:real_stopping_rule}
\end{equation}
Если ни одно условие не выполнено, полагается \(\tau_t=K\).

Ниже оставлены пять правил. Они выбраны потому, что легко реализуются, имеют ясную интерпретацию и дают разные режимы управления стоимостью.

\subsection{Правило накопленной вероятности}

Первое правило останавливается, когда вероятность того, что оптимальный момент уже наступил, становится достаточно высокой:
\begin{equation}
  \sum_{j\le k}\widehat p_{t,k}(j)\ge c.
  \label{eq:rule_cdf}
\end{equation}
Порог \(c\in(0,1)\) подбирается на проверочной выборке. Большой порог делает правило осторожным и увеличивает среднее число шагов. Малый порог дает более раннюю остановку.

Это правило является основным кандидатом для практического использования. Оно устойчиво к шуму в одной точке распределения и прямо учитывает случай, когда лучший момент уже прошел.

\subsection{Правило малой вероятности будущего улучшения}

Второе правило останавливается, когда вероятность будущего оптимума мала:
\begin{equation}
  \sum_{j>k}\widehat p_{t,k}(j)\le \varepsilon.
  \label{eq:rule_future}
\end{equation}
Параметр \(\varepsilon\) задает допустимую вероятность того, что продолжение еще могло бы быть полезным. В отличие от предыдущего правила, здесь внимание направлено не на накопленную массу прошлого и настоящего, а на остаточную массу будущего.

\subsection{Правило условной интенсивности}

Третье правило использует дискретную условную «интенсивность» (hazard), построенную по предсказанной мере \(\widehat p_{t,k}\), а не по неизвестному истинному закону:
\begin{equation}
  h_{t,k}
  =\frac{\widehat p_{t,k}(k)}{\sum_{j\ge k}\widehat p_{t,k}(j)}.
  \label{eq:hazard_rule}
\end{equation}
Она совпадает с классической формулой \(\Prob(\tau_t^\star=k\mid \tau_t^\star\ge k,\Hh_{t,k})\) в том специальном случае, когда \(\widehat p_{t,k}\) совпадает с истинным условным распределением.
Остановка выполняется при
\begin{equation}
  h_{t,k}\ge c_h.
  \label{eq:hazard_threshold}
\end{equation}
Это правило связывает распределительную постановку с классическим видом последовательного решения «остановиться или продолжить». Оно полезно как интерпретируемая производная от полной модели распределения.

\subsection{Правило квантиля}

Пусть \(q_\alpha(t,k)\) --- \(\alpha\)-квантиль предсказанного распределения:
\begin{equation}
  q_\alpha(t,k)=\min\left\{r:\sum_{j\le r}\widehat p_{t,k}(j)\ge \alpha\right\}.
  \label{eq:quantile}
\end{equation}
Правило останавливается, если
\begin{equation}
  q_\alpha(t,k)\le k.
  \label{eq:quantile_rule}
\end{equation}
Квантильное правило менее чувствительно к дальнему хвосту, чем правило по математическому ожиданию. Поэтому оно удобно в экспериментах с тяжелохвостыми моделями.

\subsection{Бюджетное правило}

Пятое правило не задает новый функционал, а выбирает порог для одного из предыдущих правил так, чтобы выполнялось ограничение на среднее число шагов:
\begin{equation}
  \E\tau_t\le C.
  \label{eq:budget_constraint}
\end{equation}
Например, можно использовать правило накопленной вероятности и подобрать \(c\) на проверочной выборке. Это правило необходимо для честного сравнения: две модели следует сравнивать не только при одном пороге, но и при одинаковой вычислительной стоимости.

\begin{table}[H]
\centering
\caption{Пять правил остановки}
\label{tab:stopping_rules}
\begin{tabularx}{\textwidth}{lYY}
\toprule
Правило & Условие остановки & Что контролирует \\
\midrule
Накопленная вероятность & \(\sum_{j\le k}\widehat p_k(j)\ge c\) & вероятность, что оптимум уже наступил \\
Малое будущее улучшение & \(\sum_{j>k}\widehat p_k(j)\le\varepsilon\) & остаточную пользу продолжения \\
Условная интенсивность & \(\widehat p_k(k)/\sum_{j\ge k}\widehat p_k(j)\ge c_h\) & текущий шаг среди не прошедших \\
Квантиль & \(q_\alpha(k)\le k\) & надежный уровень остановки \\
Бюджет & \(\E\tau\le C\) & среднюю стоимость вычислений \\
\bottomrule
\end{tabularx}
\end{table}

